\subsection{Teaching Algorithms}

The literature generally recommends making algorithms more concrete, contextualized, and interactive to support meaningful student understanding of
  algorithms.
Gal-Ezer and Zur found that high school students frequently develop misconceptions about algorithmic efficiency, often assuming shorter programs or
  programs using fewer variables are automatically more efficient~\cite{gal2004efficiency}.
These misconceptions suggest that students often rely on intuitive reasoning rather than formal computational analysis when evaluating algorithms.
Similarly, research examining secondary students’ understanding of algorithms found that many learners struggle with ideas such as testing,
  debugging, algorithm representation, and identifying meaningful edge cases~\cite{basu2023exploring}.

Because of these cognitive difficulties, Oshanova et al.\ argue that algorithmic thinking develops most effectively when instruction progresses from
  practical and familiar examples toward more abstract computational structures~\cite{Oshanova2019}.
Their work emphasizes the importance of educational “executors” such as robots, turtles, and other visual agents.
Additionally, Thorgeirsson et al.\ investigated the teaching of recursion through visual programming environments and found that such tools
  significantly improved students’ conceptual understanding of recursive processes compared to more traditional block-based programming
  approaches~\cite{Lais2024}.
The researchers argue that visual representations help students trace control flow, understand function calls, and conceptualize recursive
  decomposition, which are often difficult for new students.

Nijenhuis-Voogt et al.\ examined teachers’ perspectives on using real-world contexts to teach algorithms in secondary education and found that
  authentic, relatable contexts can improve student engagement and understanding~\cite{Voog2012}.
Their findings suggest that students benefit when algorithmic concepts are connected to familiar or meaningful scenarios, such as social media,
  games, transport, or problems perceived by the student in their lives.
Context-based learning is believed to reduce the perceived mental-load of algorithms by relating computational thinking to recognizable experiences.
However, the study also noted that teachers face challenges selecting contexts that are simultaneously authentic, motivating, and pedagogically
  useful for diverse student populations.

Romeike emphasizes the importance of creativity and engagement in algorithm instruction and argues that educators often underutilize creativity
  despite the inherently creative nature of algorithm design and problem solving~\cite{rome2007}.
Creative, open-ended programming tasks were shown to increase student motivation, improve attitudes toward computer science, and encourage
  experimentation and exploration.
Allowing students to create games, animations, or projects that the student can personally connect with helped students connect algorithmic concepts
  to self-expression and allow for authentic problem solving~\cite{Voog2012}.

Literature additionally highlights the importance of questioning, scaffolding, and guided questions during algorithm instruction.
Kátai and Osztián examined the use of dance-based algorithm visualizations combined with teacher-guided questioning to support computational
  thinking~\cite{Zoltan2021}.
Their findings suggest that carefully designed questioning strategies can help students identify algorithmic patterns, compare efficiency, and reason
  about best-case and worst-case behavior even when students have little prior programming experience.
These findings reinforce constructivist perspective, such as Romeike’s work, where students develop understanding through guided exploration rather
  than passive reception of information.

There are several disagreements that emerge regarding the most effective instructional approaches, such as the use of contextualized learning versus
  abstract or context-free instruction.
Research completed by Nijenhuis-Voogt et al.\ acknowledges that her studies run counter to that of some researchers advocating teaching programming
  concepts in isolation, arguing that students require extensive direct practice with abstract concepts before meaningful transfer can
  occur~\cite{Voog2012}.
Though there are disagreements within the presented literature concerning the complexity of topics appropriate for secondary students.
Gal-Ezer and Zur work on deeper understanding is at odds with numerous studies documenting that secondary students already struggle with basic
  abstraction, sequencing, debugging, and algorithm representation~\cite{gal2004efficiency}.

Studies also disagree on the required balance between creativity and structured scaffolding.
Romeike argues that creativity, experimentation, and open-ended project work are essential, while Oshanova et al.\ advocates gradual progression from
  concrete examples toward more abstract algorithmic~\cite{Voog2012}.
There are several studies that also conclude that algorithm visualizations significantly improve conceptual understanding by helping students observe
  control flow and algorithm execution directly.
However, Kátai and Osztián argue that this alone is insufficient for developing deeper conceptual understanding and that teacher-guided questioning
  and discussion are necessary for students to recognize core concepts~\cite{Zoltan2021}.
